\begin{prob}
    Consider the code below:

    \begin{minted}[autogobble]{python}
        def foo(n):
            i = 1
            while i * i < n:
                i += 1
            return i
    \end{minted}

    \begin{subprobset}
        \begin{subprob}
            What does \python{foo(n)} compute, roughly speaking?

            \begin{soln}
                It computes, approximately, $\sqrt{n}$. Of course, \python{foo}
                always returns an integer, so the result of the function is usually
                not exactly correct.

                More precisely, \python{foo} returns the largest integer greater
                than or equal to $\sqrt n$. For example, $\sqrt 5$ is between 2 and 3;
                \python{foo(5)} returns 3.
            \end{soln}
        \end{subprob}

        \begin{subprob}
            What is the asymptotic time complexity of \python{foo}?

            \begin{soln}
                $\Theta(\sqrt n)$
            \end{soln}
        \end{subprob}
    \end{subprobset}

\end{prob}
